\documentclass[../../../include/open-logic-section]{subfiles}

\begin{document}

\olfileid{sfr}{infinite}{card-sb}

\olsection{Bijlage: bewijs van de stelling van Schr\"oder-Bernstein}

Voordat we de na\"ieve verzamelingenleer verlaten, behandelen we nog één
laatste na\"ief (maar verfijnd!) bewijs. Het betreft een bewijs van de
stelling van Schr\"oder-Bernstein
(\olref[sfr][siz][sb]{thm:schroder-bernstein}): als
$\cardle{A}{B}$ en $\cardle{B}{A}$, dan $\cardeq{A}{B}$; met andere
woorden, bij injecties $f \colon A
\to B$ en $g \colon B \to A$ bestaat er een bijectie $h \colon A
\to B$.

In dit hoofdstuk hebben we Dedekinds begrip van \emph{afsluitingen}
gebruikt. Dedekind gaf met dit begrip zelfs een fraai bewijs van de
stelling van Schr\"oder-Bernstein; dat bewijs geven we hier. Onze
weergave volgt \citet[pp.~157--8]{Potter2004} op de voet. Daar is een
enigszins andere, maar in wezen vergelijkbare behandeling te vinden.
Een korte zoekopdracht op internet laat bovendien zien dat voor deze
stelling---net als voor de irrationaliteit van $\sqrt{2}$---\emph{veel}
interessante, uiteenlopende bewijzen bestaan.

Naar analogie met de notatie in
\olref[sfr][infinite][dedekind]{Closure} stellen we
\[
\Closureofunder{f}{B} = \bigcap \Setabs{X}{B \subseteq X
\text{ en $X$ is $f$-gesloten}}
\]
voor elke verzameling~$B$ en functie~$f$. De zo gedefinieerde
$\Closureofunder{f}{B}$ is de kleinste $f$-gesloten verzameling die~$B$
bevat, in de volgende zin:

\begin{lem}\ollabel{Closureprops}
	Voor elke functie $f$ en elke $B$ geldt:
	\begin{enumerate}
		\item\ollabel{Closurehaselem} $B \subseteq \Closureofunder{f}{B}$; en
		\item\ollabel{Closureclosed} $\Closureofunder{f}{B}$ is $f$-gesloten; en
		\item\ollabel{Closuresmallest} als $X$ $f$-gesloten is en $B
		\subseteq X$, dan $\Closureofunder{f}{B} \subseteq X$.
	\end{enumerate}
\end{lem}

\begin{proof}
Precies zoals in \olref[sfr][infinite][dedekind]{closureproperties}.
\end{proof}

Voor de stelling van Schr\"oder-Bernstein hebben we nog één laatste feit
nodig:

\begin{prop}\ollabel{sbhelper}
Als $A \subseteq B \subseteq C$ en $A \approx C$, dan $\cardeq{\cardeq{A}{B}}{C}$.
\end{prop}

\begin{proof}
Zij een bijectie~$f \colon C \to A$ gegeven. Stel $F =
\Closureofunder{f}{C \setminus B}$ en definieer als volgt een functie
$g$ met domein~$C$:
\[
	g(x) =
	\begin{cases}
			f(x) &\text{als $x \in F$}\\
			x & \text{anders}
		\end{cases}
\]
We zullen aantonen dat $g$ een bijectie van $C \to B$ is. Daaruit
volgt dat $\comp{f^{-1}}{g} \colon A \to B$ een bijectie is, waarmee
het bewijs is voltooid.

We beweren eerst: als $x \in F$ maar $y\notin F$, dan $g(x) \neq
g(y)$. Stel ter tegenspraak dat dit niet zo is. Dan geldt $y = g(y) = g(x) =
f(x)$. Omdat $x \in F$ en $F$ volgens \olref{Closureprops}
$f$-gesloten is, volgt $y = f(x) \in F$, een tegenspraak.

Stel nu dat $g(x) = g(y)$. Uit het voorgaande volgt dat $x \in F$ dan
en slechts dan als $y \in F$. Als $x, y \in F$, dan is $f(x) = g (x)
= g(y) = f(y)$ en dus $x = y$, omdat $f$ een bijectie is. Als
$x, y \notin F$, dan $x = g(x) = g(y) = y$. Bijgevolg is $g$ een
injectie.

Rest nog aan te tonen dat $\ran{g} = B$. Eerst bewijzen we dat
$\ran{g} \subseteq B$. Neem $z \in C$. Als $z \in F$, dan geldt
$g(z) = f(z) \in A \subseteq B$. Als $z \notin F$, dan is
$z \notin C \setminus B$, want $C \setminus B \subseteq F$. Dus is
$z \in B$ en geldt $g(z) = z \in B$. Bijgevolg is
$\ran{g} \subseteq B$. Voor de omgekeerde inclusie kiezen we nu
$x \in B \subseteq C$. Als
$x \notin F$, dan $g(x) = x$. Als $x \in F$, dan is $x = f(y)$ voor
een $y \in F$. Anders zou $F \setminus \{x\}$ namelijk $f$-gesloten
zijn en $C\setminus B$ omvatten, hetgeen volgens \olref{Closureprops}
onmogelijk is. Nu geldt $g(y) = f(y) = x$.
\end{proof}

Ten slotte volgt het bewijs van de hoofdstelling. Bedenk dat we voor een
functie~$h$ en verzameling~$D$ definiëren: $\funimage{h}{D} =
\Setabs{h(x)}{x \in D}$.

\begin{proof}[Bewijs van de stelling van Schr\"oder-Bernstein] Zij $f
\colon A \to B$ en $g \colon B \to A$ injecties. Omdat
$\funimage{f}{A} \subseteq B$, geldt
$\funimage{g}{\funimage{f}{A}} \subseteq g[B] \subseteq A$. Ook is
$\comp{f}{g} \colon A \to \funimage{g}{\funimage{f}{A}}$ een
injectie, omdat zowel $g$ als $f$ dat is. Volgens de definitie van zijn
codomein is $\comp{f}{g}$ zelfs een bijectie. Dus
$\cardeq{\funimage{g}{\funimage{f}{A}}}{A}$, en uit
\olref{sbhelper} volgt dat er een bijectie $h \colon A \to
\funimage{g}{B}$ bestaat. Bovendien is $g^{-1}$ een bijectie van
$\funimage{g}{B} \to B$. Bijgevolg is $\comp{h}{g^{-1}} \colon A \to
B$ een bijectie.
\end{proof}
\end{document}